% --------------------------------------------------------------
% This is all preamble stuff that you don't have to worry about.
% Head down to where it says "Start here"
% --------------------------------------------------------------

\documentclass[12pt]{article}

\usepackage[margin=1in]{geometry} 
\usepackage{amsmath,amsthm,amssymb}
\usepackage{multicol} 
\setlength{\columnsep}{1cm} % Adjust space between columns
\setlength{\parindent}{0pt} 

\everymath{\displaystyle}
\newcommand{\N}{\mathbb{N}}
\newcommand{\Z}{\mathbb{Z}}

\newenvironment{theorem}[2][Theorem]{\begin{trivlist}
\item[\hskip \labelsep {\bfseries #1}\hskip \labelsep {\bfseries #2.}]}{\end{trivlist}}
\newenvironment{lemma}[2][Lemma]{\begin{trivlist}
\item[\hskip \labelsep {\bfseries #1}\hskip \labelsep {\bfseries #2.}]}{\end{trivlist}}
\newenvironment{exercise}[2][Exercise]{\begin{trivlist}
\item[\hskip \labelsep {\bfseries #1}\hskip \labelsep {\bfseries #2.}]}{\end{trivlist}}
\newenvironment{problem}[2][Problem]{\begin{trivlist}
\item[\hskip \labelsep {\bfseries #1}\hskip \labelsep {\bfseries #2.}]}{\end{trivlist}}
\newenvironment{question}[2][Question]{\begin{trivlist}
\item[\hskip \labelsep {\bfseries #1}\hskip \labelsep {\bfseries #2.}]}{\end{trivlist}}
\newenvironment{corollary}[2][Corollary]{\begin{trivlist}
\item[\hskip \labelsep {\bfseries #1}\hskip \labelsep {\bfseries #2.}]}{\end{trivlist}}

\newenvironment{solution}{\begin{proof}[Solution]}{\end{proof}}

\begin{document}

% --------------------------------------------------------------
%                         Start here
% --------------------------------------------------------------

\title{MATH 116-04 \\ SI Session 1}
\author{Integration by Substitution and L'Hôpital}

\maketitle

\subsubsection*{I. Rearrange the functions if needed and find a suitable u-substitution }

\begin{multicols}{2}
\begin{enumerate}
    \item $\displaystyle \int x(1-x)^{99} \, dx$
    \item $\displaystyle \int \cos^3(x) \, dx$
    \item $\displaystyle \int \sin^2(x)\cos^3(x) \, dx$
    \item $\displaystyle \int \frac{\sec^2(x)}{\tan(x)+1} \, dx$
    \item $\displaystyle \int x\sqrt{x-5} \, dx$
    \item $\displaystyle \int \frac{x^2}{\sqrt{x-1}} \, dx$
    \item $\displaystyle \int x\sin(x^2)\cos(x^2) \, dx$
    \item $\displaystyle \int \sin(x)\cos^2(x)(1-\cos^3(x))^{9} \, dx$
    \item $\displaystyle \int \frac{x^5}{(x^3-3)^{3/2}} \, dx$
    \item $\displaystyle \int \frac{dx}{1+e^x}$
    \item $\displaystyle \int \frac{x^2-1}{(x^4+3x^2+1)\arctan\left(\frac{x^2+1}{x}\right)} \, dx$

\end{enumerate}
\end{multicols}



\subsubsection*{II. Compound Interest Problem}

An account starts with \$1.00 and pays 100 percent interest per year. If the interest is credited once, at the end of the year, the value of the account at year-end will be \$2.00. 

\begin{itemize}
    \item[\quad \textbf{a.}] Write the compound interest formula.
    \item[\quad \textbf{b.}] Suppose it is credited twice: 50 percent in the first six months, then 50 percent again at the end of the year. What is the amount in the account?
    \item[\quad \textbf{c.}] Suppose it is credited \textit{n} times. Write the formula.
    \item[\quad \textbf{d.}] Evaluate the formula you obtained for \textbf{c} as \textit{n} approaches infinity.
\end{itemize}

\vspace{5mm}
\pagebreak 
\subsubsection*{III. Evaluate the Limits (1-5)}

\begin{multicols}{2}
\begin{enumerate}
    \item $\displaystyle \lim_{x \to \infty} \left(1 + \frac{1}{x}\right)^{7x}$ 
    \item $\displaystyle \lim_{x \to \infty} \left(1 + \frac{7}{x}\right)^{x}$ 
    \item $\displaystyle \lim_{x \to 0} \left(1 + 7x\right)^{\frac{1}{x}}$ 
    \item $\displaystyle \lim_{x \to \infty} \left(1 + \frac{1}{x^2}\right)^x$
    \item $\displaystyle \lim_{x \to \infty} \left(\frac{x}{1+x}\right)^x$ 
    \item $\displaystyle \lim_{x \to \infty} \left(\frac{2x^2+3}{2x^2+5}\right)^{8x^2+3}$
    \item $\displaystyle \lim_{x \to 0} \left(\frac{1+\tan x}{1+ \sin x}\right)^{1/ \sin x}$
    \item $\displaystyle \lim_{x \to 1} \left(1 + \sin(\pi x)\right)^{\cot(\pi x)}$

\end{enumerate}
\end{multicols}

\vspace{5mm}

\subsubsection*{IV. If You Have Time: Extra Problems}

\begin{itemize}
    \item[\quad \textbf{a.}] Given the limit  
    \[\displaystyle \lim_{x \to \infty} f(x)^{g(x)},\]  
    where \( f(x) \to 1 \) and \( g(x) \to \infty \), prove that this limit is given by:  
    \[\displaystyle e^{\lim_{x \to \infty} g(x) [f(x) - 1]}\]
    
    \vspace{20mm}
    \item[\quad \textbf{b.}] Does the following limit exist? \[\displaystyle \lim_{x \to \pi/2}(\sin x)^\frac{1}{\cos x}\]
    
    \vspace{20mm}
    \item[\quad \textbf{c.}] Find the antiderivatives of the function $f:[0,2]\to \mathbb{R}$,
    \[
    f(x) = \sqrt{x^3+2-2\sqrt{x^3+1}}+\sqrt{x^3+10-6\sqrt{x^3+1}}
    \]
\end{itemize}

\vspace{5mm}

% --------------------------------------------------------------
%     You don't have to mess with anything below this line.
% --------------------------------------------------------------

\end{document}
